\section{\hesp{} 设计}
\label{sec:approach}

\hesp{} 位于本地 LLM 与探针目录之间（\cref{fig:pipeline}）。模型负责提议：每一步，它可以建议一个探针、提出一个结论或弃权，并在探针返回文本时解析日志。控制器负责决定：运行哪个探针、接受哪个结论、何时结束调查。每一步都在其观测到达之前写入日志，事后接受审计。

我们用运行示例走一遍。告警留下九个候选原因，各为 0.111。控制器按单位成本的期望信息增益对探针排序，运行 \texttt{source\_ips}，尽管模型提议的是 \texttt{auth\_log}（每成本单位 1.74 比 0.51 比特）；结果为 \textit{normal}，剩下四个原因，各为 0.238。\texttt{access\_pattern} 之后剩下三个，各为 0.327：暴露的管理面板、存在漏洞的组件、DNS 命令控制。\texttt{admin\_exposure} 返回 \textit{exposed, no auth}，把管理面板抬到 0.966。在下一次调用模型之前，控制器发现领先原因已超过接受阈值、且最后一条观测能把它单独识别出来，于是在三个探针、5.1\,s 之后以这个结论结束回合。

\begin{figure*}[t]
  \centering
  \includegraphics[width=\textwidth]{fig_pipeline.pdf}
  \caption{运行示例上的 \hesp{}（Llama-3.1-8B；真实原因：暴露的管理面板）。\protect\stepnum{1} 账本维护候选原因上的后验，依据冻结的探针—结果表用贝叶斯规则更新。\protect\stepnum{2} 控制器运行单位成本期望信息增益最高的合法探针；模型的提议（这里三次都是 \texttt{auth\_log}）只作参考。\protect\stepnum{3} 一旦领先原因的后验达到 0.8、且有一条当前观测能把它单独识别出来，控制器就给出结论并引用该证据。验证器只用于评测。所有数值取自归档日志。}
  \label{fig:pipeline}
\end{figure*}

\subsection{证据账本}
\label{sec:ledger}
\label{sec:tables}
\emph{问题。} 控制器需要一个模型，告诉它每个探针在每个原因下会返回什么；LLM 给不出这样的模型（向规划器询问得到的表让解决率停在 0.069 到 0.375，\S\ref{sec:policy}）。

\emph{机制。} 部署之前，一个脚本在每个原因的 $k$ 个开发回合中把每个探针各运行一次，并把结果计数成表
\[\widehat P(o \mid h, a) = \frac{q_{h,a,o}}{\sum_{o'} q_{h,a,o'}},\quad q_{h,a,o} = (1-\varepsilon)\,\frac{n_{h,a,o}}{n_{h,a}} + \frac{\varepsilon}{|O_a|},\]
其中 $n_{h,a,o}$ 是原因 $h$ 下探针 $a$ 得到结果 $o$ 的有效观测数，$O_a$ 是探针的结果词表，$\varepsilon = 0.01$。没有开发案例的原因（如 \textit{other}）取均匀行。表在使用前计算哈希并冻结。对每条告警，账本从候选原因上的均匀先验开始，在每条有效观测后应用贝叶斯规则 $p_t(h) \propto p_{t-1}(h)\,\widehat P(o_t \mid h, a_t)$，并记录该观测提高还是降低了每个原因的分数。瞬时故障（如 HTTP 503）写入日志，但不算证据。环境报告事件变化时，分数重置为先验，此前的观测不再算作当前证据。

\emph{替代方案。} 向模型询问似然可以省去开发数据，但如上所述行不通。在 sec-triage 中，每个原因五个开发回合得到的表就已能解决每一个案例（\S\ref{sec:policy}）；部署时可以对已结案工单或重放的事件计数（\S\ref{sec:discussion}）。

\subsection{探针选择}
\label{sec:selection}
对每个合法探针（即在当前状态下未运行过、在预算之内且前置条件满足），控制器计算期望信息增益
\begin{align*}
\mathrm{EIG}(a) &= H(p) - \textstyle\sum_{o} P(o \mid a)\, H\big(p(\cdot \mid o, a)\big),\\
P(o \mid a) &= \textstyle\sum_h p(h)\, \widehat P(o \mid h, a),
\end{align*}
并运行 $\mathrm{EIG}(a)/c(a)$ 最高的探针，$c(a)$ 为其成本。这一准则是经典的~\citep{dekleer1987,golovin2010}；关键在于模型的提议不再决定运行哪个探针。模型收到任务、目录、观测、账本和被拒提议的反馈（附录~\ref{app:prompt}），返回一个 JSON 决策。

\subsection{证据规则与停止}
\label{sec:finish}
\emph{结束守卫。} 只有当所称原因的当前分数不低于 $\tau = 0.8$、且至少一条被引用的观测是当前的并在账本中支持它时，模型提出的结论才会被接受。被拒绝的结论作为反馈返回给模型。日志里的文本从不进入账本，因此单凭文本无法满足守卫。

\emph{控制器停止。} 在每次调用模型之前以及预算耗尽时，控制器对当前领先原因应用同样的规则，满足时自行给出结论，并引用至多三条支持观测。\emph{确认}变体还要求有一条观测\emph{特异地}支持领先原因，即其结果在领先原因下的概率至少是任何其他具名原因下的十倍。它只使用预测表，并禁止纯粹靠排除得出的结论。

\emph{恢复。} 模型可以通过反复提出略低于阈值的同一结论，让带守卫的控制器卡住。连续两次结论被拒后，控制器直接运行自己排名第一的探针，而不再询问模型。

\subsection{跨来源佐证}
\label{sec:corroboration}
一旦由控制器决定，攻击者就改为伪造证据，而单条特异地支持良性原因的伪造观测就足够了，因为控制器先按信息量最大的观测行动。因此，\hesp{} 要求良性结论——无论由模型提出还是由控制器停止得出——建立在至少两个不同\emph{来源组}的当前观测之上，且每条都特异地支持它。来源组是探针背后的上游数据源，按部署声明；读取同一信誉源的两个探针属于同一组。规则不满足时控制器继续探测；如果证据始终不一致，案件就被升级。

\subsection{读取器信任}
\label{sec:readertrust}
当探针返回日志文本时，由解析器把它映射到探针的结果词表。按文档格式编写的规则解析器是安全的，但格式一变就什么也读不出；模型能读任何格式，却可能被说服。\hesp{} 给每条观测标记产生它的解析器。在\emph{读取器信任}规则下，模型解析的观测会更新后验，也可以确认需处置的原因，但在控制器停止、佐证和结束守卫中，都永远不算作对良性结论的特异支持。这种不对称来自攻击者的目标：确认攻击的伪造读数只造成一次不必要的升级，而确认良性原因的伪造读数会关闭案件。这条规则是一种来源标签，与能力（capability）和信息流标签~\citep{camel,fides}同一思路，只是专门用于一个决定。

\subsection{安全分析}
\label{sec:analysis}
无论模型输出什么，以下三条性质都成立：每个被执行的探针都在只读目录内且在预算之内；每个被接受的结论都引用了一条提高所称原因分数的当前观测（G3）；每条预测都在它所预测的观测之前写入日志，由审计检查。针对 A1，结束守卫和控制器停止就足够了，因为注入的文本从不成为账本证据。针对 A2，两者都无效：伪造的来源产生的是看起来真实的证据。只要攻击者至多控制一个来源组，跨来源佐证就能恢复 G2。针对 A3，佐证不再有效，因为攻击者的文本随每个探针的响应一起到达，可以伪造每个来源组的读数。只要存在可信解析器，读取器信任就能恢复 G1，代价是只有模型能读日志时，诚实的良性案例会被升级。这些规则都不能让结论正确；正确与否取决于预测表以及真实原因是否在候选之中，\S\ref{sec:policy} 对此进行评估。

\subsection{实现}
\label{sec:implementation}
\hesp{} 约 2{,}900 行 Python，只使用标准库。同一个 JSON 决策协议服务于 vLLM 和 Ollama 端点，以及用作不含 LLM 参照的脚本规划器。每个回合在自己的本地回环 HTTP 沙箱中运行，因此探针是真实的 HTTP 请求，会遇到瞬时故障和状态变化。预写式日志记录每个请求、决策、预测、观测和账本更新，审计重放每份日志并检查事件顺序、来源、预算、引用以及每个控制器结论。201 个单元测试覆盖账本、选择器、守卫、停止、佐证、读取器信任，以及预测表估计器与环境生成器的隔离。
